\documentclass[12pt,a4paper]{article}

% Кодировки и русский язык
\usepackage{cmap}
\usepackage[T2A]{fontenc}
\usepackage[utf8]{inputenc}
\usepackage[russian]{babel}

% Стандартный шрифт Computer Modern (с кириллицей через CM-Super)
\usepackage{amsmath}
\usepackage{amssymb}

\usepackage[left=2cm,right=2cm,top=2cm,bottom=2.5cm]{geometry}
\usepackage{hyperref}
\usepackage{enumitem}
\usepackage{xcolor}
\usepackage{titlesec}
\usepackage{setspace}
\usepackage{microtype}

\setstretch{1.15}

\hypersetup{
    colorlinks=true,
    linkcolor=black,
    urlcolor=blue,
    pdftitle={Концепция проекта: Web-органайзер и оркестратор учебного процесса репетитора},
    pdfauthor={ITMO}
}

\microtypesetup{
    protrusion=true,
    expansion=true,
    nopatch=footnote
}

\titleformat{\section}{\large\bfseries}{\thesection.}{0.5em}{}
\titleformat{\subsection}{\normalsize\bfseries}{\thesubsection.}{0.5em}{}
\titleformat{\subsubsection}{\small\bfseries}{\thesubsubsection.}{0.5em}{}

\titlespacing*{\section}{0pt}{1.2em}{0.4em}
\titlespacing*{\subsection}{0pt}{0.8em}{0.3em}

\setlist{nosep, topsep=2pt, itemsep=2pt, leftmargin=1.4em}

\begin{document}

\begin{center}
    {\normalsize\textbf{Концептуальный документ (Project Vision \& Functional Scope)}}\\[0.3em]
    {\Large\bfseries «Web-органайзер и оркестратор учебного процесса репетитора»}\\[0.6em]
    \rule{\textwidth}{0.6pt}
\end{center}

\vspace{0.2em}

\section{Цель и назначение системы}

Создание специализированной веб-системы для частных преподавателей и репетиторов с высокой учебной нагрузкой. Система автоматизирует рутинные процессы: распределение временных слотов, адаптивное ведение учебных программ, подготовку к урокам и консолидацию учебных материалов в едином рабочем пространстве.

\section{Формализованное описание функциональных модулей}

\subsection{Модуль 1. Интеллектуальное расписание и адаптивный календарь}

\begin{itemize}
    \item \textbf{Проблема:} При появлении нового ученика репетитор тратит много времени на сопоставление окон, учет перерывов между уроками и ручную перестройку существующей сетки занятий.
    \item \textbf{Формальное описание:}
    \begin{itemize}
        \item \textbf{Сетка регулярных и разовых занятий:} Поддержка цикличных еженедельных слотов (например, каждый вторник и четверг в 17:00) с возможностью назначения индивидуальных разовых уроков.
        \item \textbf{Помощник подбора окон (Smart Slot Finder):} Автоматический фильтр свободных интервалов времени с настраиваемым параметром «буферное время» (отдых/переключение между занятиями, например, 10--15 минут).
        \item \textbf{Быстрый доступ к инфраструктуре урока:} Каждая карточка занятия в календаре содержит активные кнопки перехода на видеозвонок (Zoom, Яндекс Телемост, Google Meet, VK Звонки) и виртуальную доску.
    \end{itemize}
\end{itemize}

\subsection{Модуль 2. Интерактивный учебный план (Curriculum Engine)}

\begin{itemize}
    \item \textbf{Проблема:} Учебный процесс часто отклоняется от плана из-за отмен и переносов; преподавателю приходится вручную пересчитывать даты тем или вести разрозненные заметки.
    \item \textbf{Формальное описание:}
    \begin{itemize}
        \item \textbf{Мастер-шаблоны курсов:} Создание типовых образовательных программ (например, «ОГЭ по математике: геометрия», «Python с нуля») с последовательной структурой тем.
        \item \textbf{Персонализация программы:} Привязка шаблона к конкретному ученику с возможностью редактирования и автоматической синхронизация с шаблоном.
        \item \textbf{Автоматическая синхронизация со статусом урока:}
        \begin{itemize}
            \item \textbf{При отмене занятия:} статус урока меняется на «Отменен», а привязанный пункт учебного плана автоматически сдвигается на следующее запланированное занятие, сдвигая вперед всю оставшуюся цепочку тем.
            \item \textbf{При переносе занятия:} преподаватель указывает новую дату/время; тема жестко перемещается в выбранный слот, а общая последовательность тем в основном расписании остается ненарушенной.
        \end{itemize}
    \end{itemize}
\end{itemize}

\subsection{Модуль 3. Панель подготовки к урокам (Pre-Lesson Dashboard)}

\begin{itemize}
    \item \textbf{Проблема:} При большом количестве учеников в день сложно удержать в голове, для кого уже готово домашнее задание, подобраны ли задачи и проверены ли прошлые работы.
    \item \textbf{Формальное описание:}
    \begin{itemize}
        \item \textbf{Дашборд на завтра / сегодня:} Виджет, собирающий список уроков на ближайшие сутки в формате чек-листа готовности.
        \item \textbf{Критерии готовности (Checklist Items):}
        \begin{itemize}
            \item Проверено ДЗ с прошлого урока;
            \item Сформированы теоретические/практические материалы к занятию;
            \item Прикреплены ссылки на видеосвязь и доску;
            \item Сформулировано домашнее задание на следующий урок.
        \end{itemize}
        \item \textbf{Индикатор прогресса:} Процент готовности преподавателя к учебному дню (помогает наглядно видеть объем оставшейся подготовки).
    \end{itemize}
\end{itemize}

\subsection{Модуль 4. Файловое хранилище и раздаточные материалы}

\begin{itemize}
    \item \textbf{Проблема:} Материалы (файлы с кодом, PDF с вариантами, презентации) теряются в мессенджерах и на локальном диске, а дисковое пространство сервера требует контролируемого расхода.
    \item \textbf{Формальное описание:}
    \begin{itemize}
        \item \textbf{Контекстная привязка файлов:} Возможность загрузки файлов (\texttt{.pdf}, \texttt{.txt}, \texttt{.py}, \texttt{.png}, \texttt{.jpg} и др.) как к конкретной теме учебного плана, так и к определенному занятию.
        \item \textbf{Планшет урока (Lesson Workspace):} В день урока в карточке занятия автоматически открывается блок «Материалы урока», где в один клик доступны все прикрепленные документы и скрипты.
        \item \textbf{Квотирование дискового пространства:} Ограничение бесплатного хранилища на пользователя в размере \textbf{1 Гб}. В интерфейсе отображается индикатор занятого места (progress-bar); при превышении квоты выводится системное уведомление с предложением расширения лимита (тарифная заглушка / демо-подписка).
    \end{itemize}
\end{itemize}

\subsection{Модуль 5. Интеграция рабочего пространства (доски и коммуникация)}

\begin{itemize}
    \item \textbf{Проблема:} Необходимость переключаться между вкладками и искать актуальную доску конкретного ученика.
    \item \textbf{Формальное описание:}
    \begin{itemize}
        \item \textbf{Гибридный подход к доскам:}
        \begin{enumerate}
            \item Поле для постоянной внешней ссылки (Miro, sBoard, Conceptboard и др.), закрепленной за учеником.
            \item Встроенный модуль легкой интерактивной доски на базе открытого решения для быстрых схем и расчетов без сторонних аккаунтов.
        \end{enumerate}
    \end{itemize}
\end{itemize}

\section{Перспективы развития (Roadmap на будущее)}

Функционал, выносимый за рамки MVP на этапе курсовой, но формулируемый как направление масштабирования:

\begin{enumerate}
    \item \textbf{Личный кабинет ученика (Read-only / Interactive):}
    \begin{itemize}
        \item Доступ ученика по персональной ссылке или пин-коду к своему расписанию, домашним заданиям, балансу занятий и архиву материалов.
    \end{itemize}

    \item \textbf{Встроенный WebRTC-сервис видеозвонков:}
    \begin{itemize}
        \item Реализация видеоконференций внутри браузера без необходимости сторонних сервисов.
    \end{itemize}

    \item \textbf{Telegram-бот ассистент:}
    \begin{itemize}
        \item Утренние напоминания преподавателю о расписании и статусе чек-листов, а также автоматическая отправка напоминаний ученикам за 2 часа до начала занятия.
    \end{itemize}
\end{enumerate}

\end{document}
